\documentclass[11pt]{article}

\usepackage[margin=1in]{geometry}
\usepackage{amsmath,amssymb,amsthm}
\usepackage{algorithm}
\usepackage{algpseudocode}
\usepackage{booktabs}
\usepackage{natbib}
\usepackage[colorlinks=true,linkcolor=blue,citecolor=blue,urlcolor=blue]{hyperref}

\newcommand{\N}{\mathbb{N}}
\newcommand{\R}{\mathbb{R}}
\newcommand{\E}{\mathbb{E}}
\newcommand{\Poi}{\mathrm{Poi}}
\newcommand{\Gam}{\mathrm{Gam}}
\newcommand{\Cat}{\mathrm{Cat}}
\newcommand{\KL}{\mathrm{KL}}
\newcommand{\MiSo}{\operatorname{MiSo}}
\newcommand{\cone}{\operatorname{cone}}
\newcommand{\one}{\mathbf{1}}

\theoremstyle{plain}
\newtheorem{theorem}{Theorem}
\newtheorem{corollary}[theorem]{Corollary}
\theoremstyle{definition}
\newtheorem{remark}{Remark}

\title{MiSo: Mixture-of-Submanifolds Nonnegative Matrix Factorization for Interpretable Factor-Usage Motifs}
\author{Dat Do}
\date{}

\begin{document}

\maketitle

\begin{abstract}
Nonnegative matrix factorization is widely used to represent a nonnegative data
matrix through low-rank nonnegative factors and loadings. Many variants
encourage individual observations to use only a small number of factors, but
they usually do not model recurring patterns of factor usage as inferential
objects. We introduce MiSo (Mixture of Submanifolds), a probabilistic model for
count data that learns recurring factor-usage motifs. Each observation is
assigned softly to a latent motif, and each motif defines a low-dimensional
submanifold spanned by a small subset of shared factor rows. The factor rows
themselves may be dense. Each motif dimension has a posterior distribution over
the shared factor rows, giving both clustering by factor-usage motif and
posterior uncertainty over the factors defining each motif. The method combines
an initial Poisson NMF dictionary with row-wise Poisson SuSiE and a finite
mixture model. We describe the model, a variational empirical-Bayes fitting
procedure, and the simulation studies needed to evaluate factor recovery,
motif recovery, overfitting behavior, and posterior uncertainty.
\end{abstract}

\section{Introduction}

Nonnegative matrix factorization (NMF) seeks a low-rank representation of a
nonnegative matrix $Y \in \N^{N \times M}$ by writing
\[
  Y \approx L F,
\]
where $L \in \R_+^{N \times K}$ contains nonnegative loadings and
$F \in \R_+^{K \times M}$ contains nonnegative factors. In many scientific
applications the rows of $L$ are expected to be sparse: each observation should
be explained by only a few latent factors. Sparse NMF methods address this by
regularizing or constraining the loading matrix. However, sparsity itself is not
the whole inferential target. Often we want to know whether many observations
share the same kind of sparsity pattern, which factors define that pattern, and
whether the evidence distinguishes one possible factor combination from another.

This paper develops MiSo, a mixture-of-submanifolds model for nonnegative count
matrix factorization. The main idea is to treat recurring factor-usage patterns
as latent motifs. Each motif represents a small collection of factor
rows that tend to appear together in the loadings. Observations are softly
assigned to motifs, and each motif has posterior uncertainty over which factors
define its dimensions. Thus the model is designed to answer questions such as:
Which factor-usage motifs are present? Which observations use each motif? Which
factors define a motif? Are two nearly identical or correlated factors
interchangeable in the posterior?

The statistical challenge is that sparse NMF has several sources of
non-identifiability. Factor rows can be permuted, highly correlated factors can
compete to explain the same counts, and an overfitted factorization can produce
duplicate or redundant factors. A hard sparse loading estimate can hide this
ambiguity by choosing one arbitrary support. We instead represent uncertainty
with posterior probabilities over factor identities, following the spirit of
SuSiE \citep{Wang2020-cr}. The resulting posterior summaries are analogous to
posterior inclusion probabilities, but they operate at the motif level rather
than only at the observation level.

The paper has three methodological layers. Section~\ref{sec:pois-susie}
introduces Poisson SuSiE for one count vector and reviews finite mixture
models. Section~\ref{sec:miso} introduces MiSo, which pools factor-usage
patterns across observations through latent motifs. Section~\ref{sec:identifiability}
is reserved for identifiability theory. Section~\ref{sec:experiments} describes
the experiments needed.

\section{Background}
\label{sec:pois-susie}
MiSo combines two familiar constructions: a SuSiE representation with a
Poisson observation model and a finite mixture model.

\subsection{Poisson SuSiE model}
A core problem in modeling a submanifold inside a large manifold is to choose the supported dimensions.
We begin with a known nonnegative dictionary
$F \in \R_+^{K \times M}$ and a count vector $y = (y_m)_{m=1}^{M}\in \N^M$ distributed as
\begin{equation}
y_m \overset{ind.}{\sim} \Poi(\mu_m), \qquad m=1,\ldots,M,
\label{eq:pois-susie-model}
\end{equation}
where $\mu = (\mu_m)_{m=1}^{M}$ is the unknown mean vector. Assume that $\mu$ is supported only on a $D$ dimensional submanifold of a $K$ dimensional manifold generated by rows $F_1, \dots, F_K \in \R^{M}$ of $F$. We model this by a SuSiE-style representation:
\begin{equation}
\mu = \sum_{d = 1}^{D} \lambda_d F_{\gamma_d},
\end{equation}
where $\gamma_d\subset [K]$ choose factor rows and $\lambda_d\in \R_+$ is the loading size for dimension $d$. They are distributed a priori according to
\begin{equation}
  \lambda_d \overset{ind.}{\sim} \Gam(\alpha_d^0,\beta_d^0),
  \qquad
  \gamma_d \overset{iid.}{\sim} \Cat(1/K,\ldots,1/K), \quad d = 1, \dots, D.
\end{equation}
where the Gamma distribution is parameterized by shape and rate. We make inference of the submanifold of $y$ via the posterior distribution of $(\gamma_d)_{d=1}^{D}$. Following Gaussian SuSiE~\cite{Wang2020-cr}, we approximate the posterior by mean-field family
\begin{equation}
  q(\lambda,\gamma) = \prod_{d=1}^D q(\gamma_d)q(\lambda_d \mid \gamma_d),
\end{equation}
which will be found by maximizing the following Evidence Lower Bound (ELBO). By the standard Jensen inequality, the model's evidence can be lower bounded as \begin{align*}
\log p(y) & \geq \mathbb E_{q} [\log p(y)] - \mathrm{KL}(q(\gamma, \lambda) \| p(\gamma, \lambda))\\
& = \sum_{m=1}^{M}\mathbb E_{q} [\log p(y_m)] - \sum_{d=1}^{D} \mathrm{KL}(q(\gamma_d, \lambda_d) \| p(\gamma_d, \lambda_d)) \\
& = \sum_{m = 1}^{M}\mathbb E_{q} \left[\log y_m \log\left(\sum_{d=1}^D \lambda_d F_{\gamma_d m}\right) - \sum_{d=1}^D \lambda_d F_{\gamma_d m} \right] - \sum_{d=1}^{D} \mathrm{KL}(q(\gamma_d, \lambda_d) \| p(\gamma_d, \lambda_d))
\end{align*}
where second line uses the independency of elements of $y$ and of the mean-field family $q$, and the third line is up to a constant ($y_m!$) only depends on data.

\begin{equation}
  q(\gamma_d=k)=\bar{\gamma}_{dk},
  \qquad
  q(\lambda_d \mid \gamma_d=k)=\Gam(\alpha_{dk},\beta_{dk}).
\end{equation}

We also introduce allocation variables $\xi_{dm}$ satisfying
$\xi_{dm}\geq 0$ and $\sum_{d=1}^D \xi_{dm}=1$ for each coordinate $m$. They
allocate the observed count $y_m$ across the $D$ dimensions. Jensen's inequality
gives the lower bound contribution
\[
  y_m \log\left(\sum_{d=1}^D \lambda_d F_{\gamma_d m}\right)
  \geq
  \sum_{d=1}^D
  \xi_{dm}y_m
  \log\left(\frac{\lambda_d F_{\gamma_d m}}{\xi_{dm}}\right).
\]

The resulting coordinate updates have a simple form. The allocation update is
\begin{equation}
  \xi_{dm}
  \propto
  \exp\left[
    \sum_{k=1}^K \bar{\gamma}_{dk}
    \{\psi(\alpha_{dk})-\log\beta_{dk}\}
    +
    \sum_{k=1}^K \bar{\gamma}_{dk}\log F_{km}
  \right],
  \label{eq:ps-xi-update}
\end{equation}
normalized over $d=1,\ldots,D$. The Gamma variational parameters are
\begin{equation}
  \alpha_{dk}
  \leftarrow
  \alpha_d^0+\sum_{m=1}^M \xi_{dm}y_m,
  \qquad
  \beta_{dk}
  \leftarrow
  \beta_d^0+\sum_{m=1}^M F_{km}.
  \label{eq:ps-gamma-update}
\end{equation}
The posterior probability that dimension $d$ uses factor $k$ is updated by a
softmax over the terms in the lower bound depending on $\gamma_d=k$:
\begin{equation}
  \bar{\gamma}_{dk}
  \propto
  \exp\left[
    \sum_{m=1}^M \xi_{dm}y_m \log F_{km}
    -
    \left(\alpha_d^0+\sum_{m=1}^M \xi_{dm}y_m\right)
    \log\left(\beta_d^0+\sum_{m=1}^M F_{km}\right)
  \right].
  \label{eq:ps-selection-update}
\end{equation}
When the rows of $F$ are normalized, the rate term is constant across $k$, but
we keep it visible because it clarifies the role of the Gamma prior and the
scale convention.

\begin{algorithm}[t]
\caption{Poisson SuSiE with a fixed dictionary}
\label{alg:poisson-susie}
\small
\begin{algorithmic}[1]
\Require Count vector $y$, row-normalized dictionary $F$, nominal dimension
$D$, initial Gamma parameters $(\alpha_d^0,\beta_d^0)$, maximum iterations
$T$, and tolerance $\delta$.
\State Initialize every row of $\bar\gamma$ to a positive probability vector,
using a small random perturbation to break symmetry.
\State Set $\xi_{dm}=1/D$,
$\alpha_{dk}=\alpha_d^0+D^{-1}\sum_m y_m$, and
$\beta_{dk}=\beta_d^0+\sum_mF_{km}$.
\For{$t=1,\ldots,T$}
  \For{$d=1,\ldots,D$}
    \State Update all $\xi_{d'm}$ using Eq.~\eqref{eq:ps-xi-update} and
    normalize over $d'=1,\ldots,D$.
    \State Update $\alpha_{dk}$ and $\beta_{dk}$ for every $k$ using
    Eq.~\eqref{eq:ps-gamma-update}.
    \State Update $\bar\gamma_{dk}$ for every $k$ using
    Eq.~\eqref{eq:ps-selection-update} and normalize over $k$.
    \If{empirical Bayes is enabled}
      \State Update $(\alpha_d^0,\beta_d^0)$ as described below.
    \EndIf
  \EndFor
  \State Evaluate the augmented variational lower bound.
  \If{its relative improvement is less than $\delta$}
    \State \textbf{break}
  \EndIf
\EndFor
\State \Return $\bar\gamma$, $\alpha$, $\beta$, $\xi$, and
$\widetilde L_k=\sum_d\bar\gamma_{dk}\alpha_{dk}/\beta_{dk}$.
\end{algorithmic}
\end{algorithm}

\paragraph{Empirical-Bayes prior update.}

Poisson SuSiE can use empirical Bayes to update the prior mean of each loading
dimension. For dimension $d$, define the posterior mean and posterior mean
log-size
\[
  \sum_{k=1}^K \bar{\gamma}_{dk}\frac{\alpha_{dk}}{\beta_{dk}},
  \qquad
  \sum_{k=1}^K
  \bar{\gamma}_{dk}\{\psi(\alpha_{dk})-\log\beta_{dk}\}.
\]
Maximizing the expected Gamma prior gives the rate update
\[
  \beta_d^0
  =
  \frac{\alpha_d^0}
       {\sum_{k=1}^K \bar{\gamma}_{dk}\alpha_{dk}/\beta_{dk}},
\]
and $\alpha_d^0$ solves
\[
  \psi(\alpha_d^0)
  =
  \sum_{k=1}^K
  \bar{\gamma}_{dk}\{\psi(\alpha_{dk})-\log\beta_{dk}\}
  + \log\beta_d^0.
\]
This empirical-Bayes step is important when $D$ is overspecified: dimensions with
little evidence can have prior mean $\alpha_d^0/\beta_d^0$ close to zero,
which is an automatic relevance determination (ARD) mechanism for selecting the
effective number of dimensions.

\subsection{Finite mixture models}

Let $Z_i\in\{1,\ldots,S\}$ denote the component of observation $Y_i$, with
$\Pr(Z_i=s)=\pi_s$. If component $s$ has likelihood $p_s(Y_i)$, then
\[
  p(Y_i)=\sum_{s=1}^S\pi_sp_s(Y_i).
\]
Suppose $\mathcal L_{is}$ is a lower bound on $\log p_s(Y_i)$. A variational
mixture update assigns responsibility
\[
  \omega_{is}=q(Z_i=s)
  =
  \frac{\pi_s\exp(\mathcal L_{is})}
       {\sum_{r=1}^S\pi_r\exp(\mathcal L_{ir})},
  \qquad
  \pi_s\leftarrow\frac{1}{N}\sum_{i=1}^N\omega_{is}.
\]
MiSo uses a Poisson SuSiE component for each motif and shares each motif's
factor-selection variables across the observations assigned to that component.

\section{MiSo}
\label{sec:miso}

\subsection{Generative model}

MiSo models repeated factor-usage motifs. The observed data are counts
$Y \in \N^{N \times M}$. The potentially dense factor matrix
$F \in \R_+^{K \times M}$ is row-normalized as above. Each motif identifies a
low-dimensional submanifold spanned by a subset of these shared factor rows;
the model does not impose sparsity on $F$. There are $S$ latent motifs. We
initially fit each motif with a
nominal dimension $D$, but the effective dimension is allowed to be smaller and
can vary by motif.

For motif $s$ and dimension $d$, the factor identity is random:
\[
  \gamma_{sd}
  \sim
  \Cat(\rho_{sd1},\ldots,\rho_{sdK}).
\]
For observation $i$,
\[
  Z_i \sim \Cat(\pi_1,\ldots,\pi_S).
\]
Conditional on $Z_i=s$, the observation has one nonnegative loading size for
each dimension:
\[
  \lambda_{id}\mid Z_i=s
  \sim
  \Gam(\alpha_{sd}^0,\beta_{sd}^0),
  \qquad d=1,\ldots,D.
\]
The observed counts are
\[
  Y_{im}\mid Z_i=s,\lambda_i,\gamma,F
  \sim
  \Poi(\mu_{im}),
  \qquad
  \mu_{im} =
  \sum_{d=1}^D \lambda_{id}F_{\gamma_{sd}m},
  \qquad m=1,\ldots,M.
\]
Equivalently, if $L_i$ denotes the loading vector for observation $i$, then
\[
  L_i
  =
  \sum_{d=1}^D \lambda_{id} e_{\gamma_{Z_i d}},
\]
where $e_k$ is the $k$th coordinate vector in $\R^K$. Thus all observations
assigned to motif $s$ share the same motif-level factor identities
$\gamma_{s1},\ldots,\gamma_{sD}$, but each observation has its own loading
sizes.

\subsection{Variational family}

The variational approximation is
\[
  q(Z,\gamma,\lambda)
  =
  \prod_{i=1}^N q(Z_i)
  \prod_{s=1}^S\prod_{d=1}^D q(\gamma_{sd})
  \prod_{i=1}^N\prod_{s=1}^S\prod_{d=1}^D
  q(\lambda_{id}\mid Z_i=s).
\]
We write
\[
  q(Z_i=s)=\omega_{is},
  \qquad
  q(\gamma_{sd}=k)=\bar{\gamma}_{sdk},
\]
and
\[
  q(\lambda_{id}\mid Z_i=s)
  =
  \Gam(\alpha_{isd},\beta_{sd}).
\]
The rate parameter $\beta_{sd}$ does not depend on $i$ because the Poisson
factor rows are shared and row-normalized; the shape parameter
$\alpha_{isd}$ depends on the counts allocated from observation $i$.

For each candidate motif $s$, we introduce allocation variables
$\xi_{isdm}$ satisfying
\[
  \xi_{isdm}\geq 0,
  \qquad
  \sum_{d=1}^D \xi_{isdm}=1.
\]
They allocate $Y_{im}$ across the dimensions of motif $s$. The implementation
computes $\xi_{isdm}$ in feature blocks rather than storing the full
$N\times S\times D\times M$ tensor.

\subsection{Coordinate updates}

The allocation update is
\begin{equation}
  \xi_{isdm}
  \propto
  \exp\left[
    \psi(\alpha_{isd})-\log\beta_{sd}
    +
    \sum_{k=1}^K \bar{\gamma}_{sdk}\log F_{km}
  \right],
  \label{eq:miso-xi-update}
\end{equation}
normalized over $d=1,\ldots,D$. The Gamma posterior update is
\begin{equation}
\begin{aligned}
  \alpha_{isd}
  \leftarrow
  \alpha_{sd}^0+\sum_{m=1}^M \xi_{isdm}Y_{im},
  \qquad
  \beta_{sd}
  \leftarrow
  \beta_{sd}^0
  +
  \sum_{m=1}^M\sum_{k=1}^K \bar{\gamma}_{sdk}F_{km}.
\end{aligned}
  \label{eq:miso-gamma-update}
\end{equation}
The second term is equal to one when each row of $F$ is normalized, but we
write the full expression to make the model dependence explicit.

For motif dimension $(s,d)$, the posterior probability of factor $k$ is updated by
\begin{equation}
  \bar{\gamma}_{sdk}
  \propto
  \rho_{sdk}
  \exp\left[
    \sum_{i=1}^N \omega_{is}
    \left\{
      \sum_{m=1}^M \xi_{isdm}Y_{im}\log F_{km}
      -
      \frac{\alpha_{isd}}{\beta_{sd}}\sum_{m=1}^M F_{km}
    \right\}
  \right].
  \label{eq:miso-selection-update}
\end{equation}
This is the central posterior uncertainty object in MiSo. A concentrated
vector $\bar{\gamma}_{sd1},\ldots,\bar{\gamma}_{sdK}$ means that motif dimension
$(s,d)$ has a well-identified factor. A diffuse vector means that several
factor rows are plausible explanations for the same motif dimension.

The motif responsibility for observation $i$ is updated from the component
lower bound for that observation under each motif:
\begin{equation}
  \omega_{is}
  \propto
  \pi_s \exp\{\mathcal{L}_{is}\},
  \label{eq:miso-responsibility-update}
\end{equation}
where
\begin{align*}
  \mathcal{L}_{is}
  &=
  \sum_{d=1}^D\sum_{m=1}^M
  \xi_{isdm}Y_{im}
  \left\{
    \psi(\alpha_{isd})-\log\beta_{sd}
    +
    \sum_{k=1}^K \bar{\gamma}_{sdk}\log F_{km}
    -
    \log \xi_{isdm}
  \right\}
  \\
  &\quad
  -
  \sum_{d=1}^D\sum_{m=1}^M
  \frac{\alpha_{isd}}{\beta_{sd}}
  \sum_{k=1}^K \bar{\gamma}_{sdk}F_{km}
  -
  \sum_{d=1}^D
  \KL\{\Gam(\alpha_{isd},\beta_{sd})
  \,\|\,\Gam(\alpha_{sd}^0,\beta_{sd}^0)\}.
\end{align*}
The mixture weights are updated by
\begin{equation}
  \pi_s \leftarrow \frac{1}{N}\sum_{i=1}^N \omega_{is}.
  \label{eq:miso-mixture-update}
\end{equation}

\subsection{Dictionary estimation and optional refinement}

When the dictionary is unknown, we first fit Poisson NMF to obtain
$Y\approx L^{\mathrm{NMF}}F^{\mathrm{NMF}}$. We row-normalize
$F^{\mathrm{NMF}}$, absorbing each removed row scale into the corresponding
column of $L^{\mathrm{NMF}}$, and denote the resulting dictionary by
$\widehat F$. This dictionary is frozen while Poisson SuSiE is fitted
independently to the rows of $Y$ for initialization. Thus the initialization
stage never alternates between the SuSiE posteriors and the factor dictionary.

The main MiSo model can either retain $F=\widehat F$ or refine it after MiSo
has been initialized. If refinement is enabled, the expected allocated counts
are
\[
  A_{km}
  =
  \sum_{i=1}^N\sum_{s=1}^S\sum_{d=1}^D
  \omega_{is}\bar{\gamma}_{sdk}\xi_{isdm}Y_{im}.
\]
The row-normalized update is
\begin{equation}
  F_{km}
  \leftarrow
  \frac{A_{km}+\epsilon}
       {\sum_{m'=1}^M(A_{km'}+\epsilon)}.
  \label{eq:miso-F-update}
\end{equation}
Keeping $F$ fixed isolates motif learning from dictionary estimation; updating
$F$ allows the dictionary to adapt to the pooled motif structure.

\subsection{Initialization}
Initialization is an important part given the non-concave nature of the optimization problem. There are many ways to initalize parameters in MiSo, depending on VI blocks. 
\paragraph{My attempt.}
With the Poisson NMF dictionary $\widehat F$ frozen, we apply
Algorithm~\ref{alg:poisson-susie} independently to each row $Y_i$. Let
$\bar\gamma^{\mathrm{PS}}_{idk}$, $\alpha^{\mathrm{PS}}_{idk}$, and
$\beta^{\mathrm{PS}}_{idk}$ denote the resulting row-wise Poisson SuSiE
parameters. Our goal is to initialize $\omega_{is}$, $\pi_s$,
$\bar\gamma_{sdk}$, $\alpha^0_{sd}$, and $\beta^0_{sd}$ from these. After
these quantities have been initialized, construct the local quantities
internally:
\[
\alpha_{isd}
=
\alpha^{0}_{sd}+\frac{\sum_mY_{im}}{D},
\]
\[
\beta_{sd}
=
\beta^{0}_{sd}
+
\sum_k\bar\gamma_{sdk}\sum_m \widehat F_{km},
\]
followed by
\[
\xi_{isdm}
\propto
\exp\left\{
\mathbb E\log\lambda_{isd}
+
\sum_k\bar\gamma_{sdk}\log \widehat F_{km}
\right\}.
\]

We aim to create $\mathcal A_i \subset [K]$ that captures the major factors
supporting observation $i$, with $|\mathcal A_i|\leq D$:
\begin{enumerate}
\item Set
\begin{equation}
  \mathrm{PIP}_{ik}
  =1-\prod_{d=1}^{D}(1-\bar\gamma^{\mathrm{PS}}_{idk}),
  \qquad k\in[K].
  \label{eq:rowwise-ps-pip}
\end{equation}
\item Let $b_{i1},\ldots,b_{iK}$ order these PIPs from largest to smallest.
For $\tau=\texttt{init\_pip\_truncate\_level}$, define
\[
  r_i
  =
  \min\left\{
    r:\frac{\sum_{j=1}^r\mathrm{PIP}_{i,b_{ij}}}
             {\sum_{k=1}^K\mathrm{PIP}_{ik}}
       \geq\tau
  \right\}.
\]
We use $\tau=0.95$ by default.
\item Set
\[
  \mathcal A_i
  =\{b_{i1},\ldots,b_{i,\min(r_i,D)}\}.
\]
Thus $\mathcal A_i$ contains at most $D$ factors, but it may contain fewer.
\end{enumerate}
Let $S^*$ be the number of unique sets among
$\{\mathcal A_i: i\in [N]\}$, let
$\mathcal A^*_{1}, \dots, \mathcal A^*_{S^*}$ be those distinct sets, and
define
\[
  \mathcal N_s^*=\{i:\mathcal A_i=\mathcal A_s^*\}.
\]
\begin{enumerate}
\item \textbf{Let's test this first.} If $S$ is not specified, set
$S=S^*$. Thus, for fixed $D$, the user only needs to specify
\texttt{init\_pip\_truncate\_level}. For $i\in[N]$ and $s\in[S]$,
initialize the responsibilities and mixture weights by
\begin{equation}
  \omega_{is}^{(0)}
  =\one\{i\in\mathcal N_s^*\},
  \qquad
  \pi_s^{(0)}=\frac{|\mathcal N_s^*|}{N}.
  \label{eq:pip-auto-omega-init}
\end{equation}

Within cluster $s$, average the row-wise PIPs,
\begin{equation}
  \overline{\mathrm{PIP}}_{sk}
  =
  \frac{1}{|\mathcal N_s^*|}
  \sum_{i\in\mathcal N_s^*}\mathrm{PIP}_{ik},
  \qquad k\in[K].
  \label{eq:cluster-average-pip}
\end{equation}
Let $a_{s1},\ldots,a_{sK}$ be a permutation of $1,\ldots,K$ such that
\[
  \overline{\mathrm{PIP}}_{s,a_{s1}}
  \geq \cdots \geq
  \overline{\mathrm{PIP}}_{s,a_{sK}}.
\]
Assuming $D\leq K$, use the first $D$ factors as distinct anchors and set
\begin{equation}
  \bar\gamma_{sdk}^{(0)}
  =
  (1-\varepsilon_\gamma)\one\{k=a_{sd}\}
  +\frac{\varepsilon_\gamma}{K},
  \qquad d\in[D],\quad k\in[K].
  \label{eq:pip-anchor-gamma-init}
\end{equation}
We use $\varepsilon_\gamma=0.05$ by default. This expands every cluster to
the common nominal dimension $D$, while the ordering of all $D$ slots is
determined by the average PIP within that cluster.

The Poisson SuSiE slot labels are exchangeable across observations. Therefore,
before initializing the Gamma priors, align the row-wise slots to the ordered
cluster anchors. For each $i\in\mathcal N_s^*$, choose
\begin{equation}
  p_i
  =
  \arg\max_{p\in\mathfrak S_D}
  \sum_{d=1}^D
  \bar\gamma^{\mathrm{PS}}_{i,p(d),a_{sd}}
  \frac{\alpha^{\mathrm{PS}}_{i,p(d),a_{sd}}}
       {\beta^{\mathrm{PS}}_{i,p(d),a_{sd}}},
  \label{eq:pip-anchor-slot-alignment}
\end{equation}
where $\mathfrak S_D$ is the set of permutations of $1,\ldots,D$. This is a
linear assignment problem, so it can be solved without enumerating the $D!$
permutations.

Directly averaging Gamma shapes and rates is generally not appropriate:
different shape--rate pairs are on a nonlinear scale and their arithmetic
averages need not represent the average posterior distribution. Instead,
average the two sufficient posterior moments after slot alignment:
\begin{align}
  \overline\lambda^{\mathrm{PS}}_{sd}
  &=
  \frac{1}{|\mathcal N_s^*|}
  \sum_{i\in\mathcal N_s^*}\sum_{k=1}^K
  \bar\gamma^{\mathrm{PS}}_{i,p_i(d),k}
  \frac{\alpha^{\mathrm{PS}}_{i,p_i(d),k}}
       {\beta^{\mathrm{PS}}_{i,p_i(d),k}},
  \label{eq:pip-cluster-mean-lambda}\\
  \overline{\log\lambda}^{\mathrm{PS}}_{sd}
  &=
  \frac{1}{|\mathcal N_s^*|}
  \sum_{i\in\mathcal N_s^*}\sum_{k=1}^K
  \bar\gamma^{\mathrm{PS}}_{i,p_i(d),k}
  \left\{
    \psi\!\left(\alpha^{\mathrm{PS}}_{i,p_i(d),k}\right)
    -\log\beta^{\mathrm{PS}}_{i,p_i(d),k}
  \right\}.
  \label{eq:pip-cluster-mean-log-lambda}
\end{align}
Because $\widehat F$ is row-normalized, the Poisson SuSiE updates make
$\alpha^{\mathrm{PS}}_{idk}$ and $\beta^{\mathrm{PS}}_{idk}$ constant in $k$
for fixed $(i,d)$. In that case, the sums over $k$ above simplify to the
posterior mean and posterior mean log-size of the aligned slot,
$\alpha^{\mathrm{PS}}_{i,p_i(d)}/\beta^{\mathrm{PS}}_{i,p_i(d)}$ and
$\psi(\alpha^{\mathrm{PS}}_{i,p_i(d)})-
\log\beta^{\mathrm{PS}}_{i,p_i(d)}$, respectively.
Initialize $\alpha^0_{sd}$ and $\beta^0_{sd}$ by matching these moments to a
Gamma distribution:
\begin{align}
  \beta^0_{sd}
  &=
  \frac{\alpha^0_{sd}}{\overline\lambda^{\mathrm{PS}}_{sd}},
  \label{eq:pip-prior-rate-init}\\
  \psi(\alpha^0_{sd})-\log\beta^0_{sd}
  &=
  \overline{\log\lambda}^{\mathrm{PS}}_{sd}.
  \label{eq:pip-prior-shape-init}
\end{align}
Equivalently, solve the one-dimensional equation
\[
  \log\alpha^0_{sd}-\psi(\alpha^0_{sd})
  =
  \log\overline\lambda^{\mathrm{PS}}_{sd}
  -\overline{\log\lambda}^{\mathrm{PS}}_{sd},
\]
and then use Eq.~\eqref{eq:pip-prior-rate-init}. A small positive numerical
floor should be applied to $\overline\lambda^{\mathrm{PS}}_{sd}$ in the
implementation. These equations use every row-wise factor possibility rather
than only the anchor, while the alignment in
Eq.~\eqref{eq:pip-anchor-slot-alignment} ensures that the same motif dimension
is averaged across observations. This is precisely one MiSo empirical-Bayes
prior update with hard initial cluster responsibilities and the row-wise
Poisson SuSiE posteriors substituted for the MiSo local posteriors.

\item If $S$ is specified and $S^* \geq S$, let
$\mathcal N^*_s=\{i:\mathcal A_i=\mathcal A^*_{s}\}$. Reorder the distinct
sets so that $(|\mathcal N^*_s|)_{s=1}^{S^*}$ is decreasing, and retain the
first $S$ sets. For $s\in[S]$, initialize the factor-selection probabilities by
\begin{equation}
  \bar\gamma_{sdk}^{(0)}
  =
  \frac{1}{|\mathcal N_s^*|}
  \sum_{i\in\mathcal N_s^*}\bar\gamma_{idk}^{\mathrm{PS}},
  \qquad d\in[D],\quad k\in[K].
  \label{eq:pip-set-gamma-init}
\end{equation}
For $i\in[N]$ and $s\in[S]$, initialize the responsibilities by
\begin{equation}
  \omega_{is}^{(0)}
  =
  \begin{cases}
    \one\{\mathcal A_i=\mathcal A_s^*\},
      & \mathcal A_i\in\{\mathcal A_1^*,\ldots,\mathcal A_S^*\},\\[3pt]
    1/S,
      & \mathcal A_i\in
        \{\mathcal A_{S+1}^*,\ldots,\mathcal A_{S^*}^*\}.
  \end{cases}
  \label{eq:pip-set-omega-init}
\end{equation}
Thus observations belonging to one of the retained sets receive a hard initial
assignment, whereas observations belonging to a discarded set are initially
distributed uniformly across the retained submanifolds.
\item If $S$ is specified and $S > S^*$: Just set $S = S^*$ and proceed. We will go back to this if the current solution is not good enough.
\end{enumerate}




\paragraph{Codex's suggestion.}

With the Poisson NMF dictionary $\widehat F$ frozen, we apply
Algorithm~\ref{alg:poisson-susie} independently to each row $Y_i$. Let
$\bar\gamma^{\mathrm{PS}}_{idk}$, $\alpha^{\mathrm{PS}}_{idk}$, and
$\beta^{\mathrm{PS}}_{idk}$ denote the resulting row-wise Poisson SuSiE
parameters. The preliminary expected loading of factor $k$ in observation $i$
is
\[
  \widetilde L_{ik}
  =
  \sum_{d=1}^D
  \bar\gamma^{\mathrm{PS}}_{idk}
  \frac{\alpha^{\mathrm{PS}}_{idk}}
       {\beta^{\mathrm{PS}}_{idk}}.
\]
We remove observation-level scale before clustering,
\[
  x_{ik}
  =
  \frac{\widetilde L_{ik}}
       {\sum_{h=1}^K\widetilde L_{ih}},
  \qquad x_i\in\Delta^{K-1},
\]
and obtain $S$ preliminary centers by solving
\[
  (\widehat z,\widehat c)
  =
  \arg\min_{z_1,\ldots,z_N,\,c_1,\ldots,c_S}
  \sum_{i=1}^N \lVert x_i-c_{z_i}\rVert_2^2,
  \qquad c_s\in\Delta^{K-1}.
\]
For preliminary cluster $s$, let
\[
  r_s(1),\ldots,r_s(K)
\]
order the factors so that
\[
  c_{s,r_s(1)}\geq c_{s,r_s(2)}\geq\cdots\geq c_{s,r_s(K)}.
\]
Assuming $D\leq K$, MiSo uses the distinct top-$D$ anchors
\begin{equation}
  \kappa_{sd}^{(0)}=r_s(d),
  \qquad d=1,\ldots,D,
  \label{eq:distinct-top-D-init}
\end{equation}
to define $D$ alignment anchors. The Poisson SuSiE slot labels are
exchangeable across observations, so they must be aligned before their
posteriors can be averaged. For each observation with $\widehat z_i=s$, choose
the permutation
\begin{equation}
  p_i
  =
  \arg\max_{p\in\mathfrak S_D}
  \sum_{d=1}^D
  \bar\gamma^{\mathrm{PS}}_{i,p(d),\kappa_{sd}^{(0)}}
  \frac{
    \alpha^{\mathrm{PS}}_{i,p(d),\kappa_{sd}^{(0)}}
  }{
    \beta^{\mathrm{PS}}_{i,p(d),\kappa_{sd}^{(0)}}
  },
  \label{eq:ps-slot-alignment}
\end{equation}
where $\mathfrak S_D$ is the set of permutations of $1,\ldots,D$. This is a
linear assignment problem. It matches each motif slot to one row-wise Poisson
SuSiE slot using the posterior expected loading assigned to the corresponding
anchor factor.

To prevent observations with large total counts from dominating the
initialization, define the within-observation loading share of Poisson SuSiE
slot $d$ by
\begin{equation}
  w_{id}
  =
  \frac{
    \sum_{k=1}^K
    \bar\gamma^{\mathrm{PS}}_{idk}
    \alpha^{\mathrm{PS}}_{idk}/\beta^{\mathrm{PS}}_{idk}
  }{
    \sum_{d'=1}^D\sum_{k=1}^K
    \bar\gamma^{\mathrm{PS}}_{id'k}
    \alpha^{\mathrm{PS}}_{id'k}/\beta^{\mathrm{PS}}_{id'k}
  }.
  \label{eq:ps-slot-weight}
\end{equation}
The cluster-averaged posterior for motif slot $(s,d)$ is then
\begin{equation}
  q_{sdk}^{(0)}
  =
  \frac{
    \sum_{i:\widehat z_i=s}
    w_{i,p_i(d)}\bar\gamma^{\mathrm{PS}}_{i,p_i(d),k}
  }{
    \sum_{i:\widehat z_i=s}w_{i,p_i(d)}
  }.
  \label{eq:aligned-ps-average}
\end{equation}
If the denominator is numerically zero, we instead set
$q_{sdk}^{(0)}=\one\{k=\kappa_{sd}^{(0)}\}$. Finally, for a small
$\varepsilon_\gamma>0$, initialize
\begin{equation}
  \bar\gamma_{sdk}^{(0)}
  =
  (1-\varepsilon_\gamma)q_{sdk}^{(0)}
  +\frac{\varepsilon_\gamma}{K},
  \qquad k=1,\ldots,K.
  \label{eq:gamma-init}
\end{equation}
We use $\varepsilon_\gamma=0.05$ by default. Unlike an initialization that
replaces every slot by a point mass on its anchor, Eq.~\eqref{eq:gamma-init}
retains the factor uncertainty learned by row-wise Poisson SuSiE. The distinct
anchors only resolve slot-label switching; they do not replace the averaged
posterior.

The initial mixture weights and responsibilities are set to
\[
  \pi_s^{(0)}=\frac{1}{S},
  \qquad
  \omega_{is}^{(0)}=\frac{1}{S}.
\]
The preliminary partition is therefore used to initialize motif factor
posteriors but is not imposed as a hard assignment. The Gamma variational and
prior parameters are initialized symmetrically across motifs and are then
updated using the raw count matrix.

\subsection{Empirical-Bayes dimension selection}

For motif $s$ and dimension $d$, define the responsibility-weighted posterior
moments
\begin{align*}
  \overline\lambda_{sd}
  &=
  \frac{
    \sum_{i=1}^N\omega_{is}\alpha_{isd}/\beta_{sd}
  }{
    \sum_{i=1}^N\omega_{is}
  },\\
  \overline{\log\lambda}_{sd}
  &=
  \frac{
    \sum_{i=1}^N\omega_{is}
    \{\psi(\alpha_{isd})-\log\beta_{sd}\}
  }{
    \sum_{i=1}^N\omega_{is}
  }.
\end{align*}
The empirical-Bayes update is
\begin{equation}
  \beta_{sd}^0=\frac{\alpha_{sd}^0}{\overline\lambda_{sd}},
  \qquad
  \psi(\alpha_{sd}^0)
  =\overline{\log\lambda}_{sd}+\log\beta_{sd}^0,
  \label{eq:miso-eb-update}
\end{equation}
where the second equation is solved numerically. Each fitted motif starts with
the same nominal dimension $D$, but its effective dimension can be smaller.
The empirical-Bayes Gamma prior mean
\[
  \frac{\alpha_{sd}^0}{\beta_{sd}^0}
\]
estimates the loading size expected for that dimension. If this value is close to
zero relative to the other dimensions in motif $s$, then dimension $d$ is effectively
inactive. This gives an ARD-style selected dimension
\[
  \widehat{D}_s
  =
  \#\left\{
    d:
    \frac{\alpha_{sd}^0/\beta_{sd}^0}
         {\sum_{d'=1}^D \alpha_{sd'}^0/\beta_{sd'}^0}
    \geq \tau
  \right\},
\]
for a threshold $\tau$.

This ARD diagnostic should be interpreted separately from the entropy of
$\bar{\gamma}_{sd1},\ldots,\bar{\gamma}_{sdK}$. The ARD prior mean asks whether
a dimension carries loading mass. The entropy of $\bar{\gamma}_{sd}$ asks whether
the factor identity of that dimension is uncertain. Both are useful: an
inactive dimension should have small loading mass, while an active but
ambiguous dimension may have large loading mass and diffuse posterior factor
identity.

\begin{algorithm}[t]
\caption{MiSo variational empirical-Bayes fitting}
\label{alg:miso}
\small
\begin{algorithmic}[1]
\Require Count matrix $Y$, factor count $K$, nominal motif dimension $D$,
PIP truncation level $\tau$, inner iterations $R$, outer iterations $T$,
tolerance $\delta$, and a choice of whether to refine $F$.
\State Fit rank-$K$ Poisson NMF; row-normalize its factor matrix to obtain
$\widehat F$ and hold $\widehat F$ fixed during the next two steps.
\For{$i=1,\ldots,N$}
  \State Apply Algorithm~\ref{alg:poisson-susie} to $Y_i$ with dictionary
  $\widehat F$; retain $\bar\gamma^{\mathrm{PS}}_{idk}$,
  $\alpha^{\mathrm{PS}}_{idk}$, and $\beta^{\mathrm{PS}}_{idk}$.
\EndFor
\State Compute Eq.~\eqref{eq:rowwise-ps-pip}; truncate each ordered PIP
profile at level $\tau$ to obtain $\mathcal A_i$, and set $S$ to the number
of distinct sets among $\mathcal A_1,\ldots,\mathcal A_N$.
\State Initialize $\omega_{is}$ and $\pi_s$ using
Eq.~\eqref{eq:pip-auto-omega-init}; order the cluster-average PIPs using
Eq.~\eqref{eq:cluster-average-pip}.
\State Initialize $\bar\gamma_{sdk}$ from the distinct top-$D$ anchors using
Eq.~\eqref{eq:pip-anchor-gamma-init}, and align the row-wise Poisson SuSiE
slots using Eq.~\eqref{eq:pip-anchor-slot-alignment}.
\State Initialize $(\alpha_{sd}^0,\beta_{sd}^0)$ by the posterior-moment
averages in Eqs.~\eqref{eq:pip-cluster-mean-lambda}--
\eqref{eq:pip-prior-shape-init}.
\State Set $F=\widehat F$, set $\xi_{isdm}=1/D$, and initialize
$\alpha_{isd}=\alpha_{sd}^0+D^{-1}\sum_mY_{im}$ and $\beta_{sd}$ using
Eq.~\eqref{eq:miso-gamma-update}.
\For{$t=1,\ldots,T$}
  \For{$r=1,\ldots,R$}
    \State Update $\xi_{isdm}$ for all $i,s,d,m$ using
    Eq.~\eqref{eq:miso-xi-update}, processing features in blocks.
    \State Update $\alpha_{isd}$ and $\beta_{sd}$ using
    Eq.~\eqref{eq:miso-gamma-update}.
  \EndFor
  \State Evaluate $\mathcal L_{is}$ for every observation and candidate motif.
  \State Update $\omega_{is}$ and $\pi_s$ using
  Eqs.~\eqref{eq:miso-responsibility-update} and
  \eqref{eq:miso-mixture-update}.
  \If{empirical Bayes is enabled}
    \State Update $(\alpha_{sd}^0,\beta_{sd}^0)$ using
    Eq.~\eqref{eq:miso-eb-update}.
  \EndIf
  \State Update $\bar\gamma_{sdk}$ using
  Eq.~\eqref{eq:miso-selection-update}; optionally damp the update.
  \If{dictionary refinement is enabled}
    \State Update $F$ using Eq.~\eqref{eq:miso-F-update}; optionally damp the
    update and row-normalize.
  \EndIf
  \State Evaluate the mixture variational objective.
  \If{its relative improvement is less than $\delta$}
    \State \textbf{break}
  \EndIf
\EndFor
\State Recompute the local posteriors and responsibilities using the final
global parameters.
\State \Return $\omega$, $\pi$, $\bar\gamma$, $\alpha$, $\beta_{sd}$,
$\widehat D_s$, and $F$.
\end{algorithmic}
\end{algorithm}

\section{Identifiability Theory}
\label{sec:identifiability}

MiSo places its geometric structure on the latent Poisson intensity, not on
the observed count vector.  This distinction lets us retain the language of
submanifolds while using cones when a closed support or an extreme ray is the
mathematically relevant object.  For a realized motif $s$, write
\begin{equation}
  H_s
  =\sum_{d=1}^{D_s}a_{sd}\delta_{\theta_{sd}},
  \qquad
  \theta_{sd}=\frac{F_{\gamma_{sd}}}{\beta_{sd}},
  \qquad a_{sd},\beta_{sd}>0,
  \label{eq:H-miso}
\end{equation}
where $a_{sd}$ and $\beta_{sd}$ are the Gamma shape and rate.  Coincident atoms
are combined by adding their shapes, giving the canonical representation
\begin{equation}
  H=\sum_{j=1}^{J}a_j\delta_{\theta_j},
  \qquad a_j>0,
  \qquad \theta_1,\ldots,\theta_J\text{ pairwise distinct}.
  \label{eq:H-canonical-main}
\end{equation}
The corresponding latent submanifold and its closure cone are
\begin{align}
  \mathcal M(H)
  &=\left\{\sum_{j=1}^{J}x_j\theta_j:x_j>0\right\},
  \label{eq:latent-submanifold}\\
  \mathcal C(H)
  &=\overline{\mathcal M(H)}
    =\cone\{\theta_1,\ldots,\theta_J\}.
  \label{eq:closure-cone}
\end{align}
Thus $\mathcal M(H)$ is the relative interior of $\mathcal C(H)$ and is a
smooth submanifold of its linear span.  Its intrinsic dimension is
$\dim\mathcal M(H)=\operatorname{rank}(\theta_1,\ldots,\theta_J)$, which can
be smaller than the nominal number of dimensions $D_s$.

To define the distribution encoded by $H$, independently draw
\[
  \xi_j\sim\Gam(a_j,1),\qquad j=1,\ldots,J,
\]
and set
\begin{equation}
  \mu_H=\sum_{j=1}^{J}\xi_j\theta_j.
  \label{eq:gamma-H-main}
\end{equation}
We denote the law of $\mu_H$ by $\Gam(H)$.  This agrees with the original
shape--rate parameterization because, if
$\lambda_j\sim\Gam(a_j,\beta_j)$, then
$\beta_j\lambda_j\sim\Gam(a_j,1)$ and
$\lambda_jF_{\gamma_j}=(\beta_j\lambda_j)(F_{\gamma_j}/\beta_j)$.
Conditional on $\mu_H$, the coordinates of $Y$ are independent with
\[
  Y_m\mid\mu_H\sim\Poi((\mu_H)_m),
  \qquad m=1,\ldots,M.
\]
Although $\mu_H\in\mathcal M(H)$ almost surely, the topological support of
its distribution is the closure cone $\mathcal C(H)$.

The population mixture of latent submanifolds is represented by
\begin{equation}
  G=\sum_{s=1}^{S}\pi_s\delta_{H_s},
  \qquad \pi_s>0,
  \qquad \sum_{s=1}^{S}\pi_s=1,
  \label{eq:G-main}
\end{equation}
and its observed count law is
\begin{equation}
  \MiSo(G)
  =\sum_{s=1}^{S}\pi_s\Poi\{\Gam(H_s)\}.
  \label{eq:MiSo-G-main}
\end{equation}
The motif measure $H_s$ need not be a probability measure: its atom masses
are Gamma shapes, whereas $\pi_s$ is the mixture probability.  For
$t\in\R_+^M$, define
\begin{equation}
  \Phi_H(t)
  :=\E\exp\{-\langle t,\mu_H\rangle\}
  =\prod_{j=1}^{J}(1+\langle\theta_j,t\rangle)^{-a_j}.
  \label{eq:Phi-H-main}
\end{equation}
The probability-generating function of the noisy counts is
\begin{equation}
  \E_H\!\left[\prod_{m=1}^{M}z_m^{Y_m}\right]
  =\Phi_H(\one-z),
  \qquad z\in[0,1]^M.
  \label{eq:pgf-main}
\end{equation}
This identity is the bridge from the observed Poisson law to the geometry of
the latent submanifolds.

\begin{theorem}[Identification of one latent submanifold]
\label{thm:H-identification-main}
Let $H$ and $H'$ be finite positive atomic measures on
$\R_+^M\setminus\{0\}$.  Then
\[
  \Poi\{\Gam(H)\}\overset{d}{=}\Poi\{\Gam(H')\}
  \qquad\Longrightarrow\qquad
  H=H'.
\]
Consequently, the latent submanifold $\mathcal M(H)$, its closure cone
$\mathcal C(H)$, and its intrinsic dimension are identifiable from the
population distribution of the Poisson observations.
\end{theorem}

\begin{proof}
Equality of the count laws gives equality of their probability-generating
functions.  By \eqref{eq:pgf-main}, $\Phi_H(t)=\Phi_{H'}(t)$ for
$t\in(0,1)^M$.  Hence their negative log-Laplace transforms agree:
\[
  \sum_j a_j\log(1+\langle\theta_j,t\rangle)
  =
  \sum_\ell a'_\ell\log(1+\langle\theta'_\ell,t\rangle).
\]
Differentiating on this open set gives
\begin{equation}
  \sum_j\frac{a_j\theta_j}{1+\langle\theta_j,t\rangle}
  =
  \sum_\ell
  \frac{a'_\ell\theta'_\ell}{1+\langle\theta'_\ell,t\rangle}.
  \label{eq:rational-main}
\end{equation}

Choose $v\in\R_{++}^M$ such that distinct vectors among the finitely many
$\theta_j$ and $\theta'_\ell$ have distinct projections onto $v$.  Such a
$v$ exists because the excluded choices lie in finitely many proper
hyperplanes.  Put $r_j=\langle\theta_j,v\rangle$ and
$r'_\ell=\langle\theta'_\ell,v\rangle$, all of which are positive.  On the
real line $t=xv$, equation \eqref{eq:rational-main} becomes
\[
  \sum_j\frac{a_j\theta_j/r_j}{x+1/r_j}
  =
  \sum_\ell\frac{a'_\ell\theta'_\ell/r'_\ell}{x+1/r'_\ell}
\]
for $x$ in a positive interval.  Multiplication by all denominators gives an
identity between vector-valued real polynomials.  Therefore the two real
partial-fraction decompositions have the same denominator roots and the same
coefficient at each root.  For a matched pair with $r_j=r'_\ell$, multiplying
by $x+1/r_j$ and taking the real limit $x\to-1/r_j$ gives
\[
  \frac{a_j\theta_j}{r_j}
  =\frac{a'_\ell\theta'_\ell}{r_j}.
\]
Taking the inner product with $v$ yields $a_j=a'_\ell$, after which
$\theta_j=\theta'_\ell$.  The generic choice of $v$ makes the matching
one-to-one, proving $H=H'$.
\end{proof}

The measure representation absorbs the benign Gamma-splitting ambiguity:
$a_1\delta_\theta+a_2\delta_\theta$ and
$(a_1+a_2)\delta_\theta$ are literally the same measure.  Theorem
\ref{thm:H-identification-main} also makes explicit that the Poisson layer is
lossless for population identification, because the count probability-
generating function reveals the Laplace transform of the latent intensity
distribution \citep{Teicher1960}.

For the mixture geometry, let
\begin{equation}
  P_G=\sum_{s=1}^{S}\pi_s\Gam(H_s)
  \label{eq:latent-mixture-main}
\end{equation}
be the distribution of the latent intensity and define
\[
  \mathcal U(G)=\operatorname{supp}(P_G),
  \qquad
  \mathcal C(G)=\cone\{\mathcal U(G)\}.
\]
The set $\mathcal U(G)$ retains the union of the closure cones of the latent
submanifolds, whereas $\mathcal C(G)$ is their global conic hull.

\begin{theorem}[Identification of support geometry and extreme rays]
\label{thm:global-cone-main}
Let $G$ and $G'$ be finite probability measures of the form
\eqref{eq:G-main}.  If
\[
  \MiSo(G)\overset{d}{=}\MiSo(G'),
\]
then
\[
  P_G=P_{G'},
  \qquad
  \mathcal U(G)=\mathcal U(G'),
  \qquad
  \mathcal C(G)=\mathcal C(G').
\]
Consequently, the extreme rays of the global closure cone are identifiable
despite the Poisson observation noise.  If rays are normalized by their
intersection with
$\Delta^{M-1}=\{x\in\R_+^M:\one^\top x=1\}$, the corresponding extreme
points are identifiable as well.  Moreover, every global extreme ray
contains an atom from at least one motif measure in the support of $G$ and
from at least one motif measure in the support of $G'$.
\end{theorem}

\begin{proof}
The count probability-generating function is the Laplace transform of $P_G$
evaluated at $\one-z$.  Equality of the MiSo laws therefore gives equality of
the multivariate Laplace transforms of $P_G$ and $P_{G'}$ on $(0,1)^M$.
Uniqueness of the Laplace transform on $\R_+^M$ gives $P_G=P_{G'}$.

For $H=\sum_j a_j\delta_{\theta_j}$, the independent Gamma coefficients in
\eqref{eq:gamma-H-main} have joint topological support $\R_+^J$.  Their
linear image consequently has support
\[
  \operatorname{supp}\{\Gam(H)\}
  =\left\{\sum_jx_j\theta_j:x_j\geq0\right\}
  =\mathcal C(H).
\]
Because the mixture is finite and every mixture weight is positive,
\begin{equation}
  \mathcal U(G)
  =\bigcup_{H\in\operatorname{supp}(G)}\mathcal C(H).
  \label{eq:union-cones-main}
\end{equation}
The same identity holds for $G'$.  Equality of $P_G$ and $P_{G'}$ therefore
identifies the union of closure cones and its global conic hull.

Extreme rays and normalized extreme points are intrinsic properties of this
identified closed cone.  Finally, $\mathcal C(G)$ is generated by the finite
union of the atoms in the component measures.  Every extreme ray of a
finitely generated cone contains a vector from any finite generating set;
otherwise a point on that ray would be a nontrivial conic combination of
vectors off the ray.  Applying this fact to the atoms under $G$ and under
$G'$ proves the final claim.
\end{proof}

When each factor row is normalized by $\one^\top F_k=1$, the scaled atom
$\theta=F_k/b$ lies on the same ray as $F_k$, and simplex normalization maps
it back to $F_k$.  Theorem~\ref{thm:global-cone-main} therefore identifies
the factor rows corresponding to global extreme rays.  It does not, by
itself, identify interior dictionary rows or assign the identified rays to
individual motifs.

\begin{remark}[Interior rays prevent unrestricted mixture identification]
\label{rem:interior-ray-main}
Theorem~\ref{thm:H-identification-main} makes the single-submanifold kernel
injective, but it does not imply that an unrestricted mixture $G$ is
identifiable.  To see this with all Gamma shapes equal to one, choose linearly
independent $\theta_A,\theta_B\in\R_+^M\setminus\{0\}$, fix
$\rho\in(0,1)$, and introduce the interior direction
\[
  \theta_C=\rho\theta_A+(1-\rho)\theta_B.
\]
Let
\[
  H_{AB}=\delta_{\theta_A}+\delta_{\theta_B},\qquad
  H_{BC}=\delta_{\theta_B}+\delta_{\theta_C},\qquad
  H_{AC}=\delta_{\theta_A}+\delta_{\theta_C}.
\]
Writing
$A(t)=1+\langle\theta_A,t\rangle$ and defining $B(t)$ and $C(t)$
analogously, we have $C=\rho A+(1-\rho)B$ and hence
\begin{align}
  \Phi_{H_{AB}}(t)
  &=\frac{1}{A(t)B(t)} \notag\\
  &=\frac{\rho}{B(t)C(t)}
    +\frac{1-\rho}{A(t)C(t)} \notag\\
  &=\rho\Phi_{H_{BC}}(t)
    +(1-\rho)\Phi_{H_{AC}}(t).
  \label{eq:interior-identity-main}
\end{align}
Thus the distinct mixing measures
\[
  G=\delta_{H_{AB}},
  \qquad
  G'=\rho\delta_{H_{BC}}+(1-\rho)\delta_{H_{AC}}
\]
produce exactly the same noisy count law.  Geometrically, one latent
submanifold has been represented by a mixture whose two closure cones
subdivide the original cone along an interior ray.  This does not contradict
Theorem~\ref{thm:global-cone-main}: both representations have the same latent
support and the same global extreme rays.  Rather, it shows why an
over-specified interior factor can prevent identification of the individual
motifs even when the global geometry is identifiable.
\end{remark}

To state the condition needed for full mixture identification, call a
transform family $\{\Phi_H:H\in\mathcal H_0\}$ finitely linearly independent
if, for every finite collection of distinct
$H_1,\ldots,H_r\in\mathcal H_0$,
\begin{equation}
  \sum_{j=1}^{r}c_j\Phi_{H_j}(t)=0
  \text{ for every }t\in(0,1)^M
  \quad\Longrightarrow\quad
  c_1=\cdots=c_r=0.
  \label{eq:finite-linear-independence-main}
\end{equation}

\begin{theorem}[Finite-mixture identification]
\label{thm:G-identification-main}
Let $G$ and $G'$ be finite probability measures on $\mathcal H_0$.  If
$\{\Phi_H:H\in\mathcal H_0\}$ is finitely linearly independent, then
\[
  \MiSo(G)\overset{d}{=}\MiSo(G')
  \qquad\Longrightarrow\qquad
  G=G'.
\]
Conversely, if the transform family is not finitely linearly independent,
then the class of finite probability mixtures on $\mathcal H_0$ is not
identifiable.
\end{theorem}

\begin{proof}
Take the union of the finite supports of $G$ and $G'$ and denote its distinct
elements by $H_1,\ldots,H_r$.  Equality of the noisy count laws and
\eqref{eq:pgf-main} gives
\[
  \sum_{j=1}^{r}\{G(H_j)-G'(H_j)\}\Phi_{H_j}(t)=0,
  \qquad t\in(0,1)^M.
\]
Finite linear independence forces every coefficient to vanish, so $G=G'$.
Conversely, consider a nonzero finite linear relation among distinct
transforms.  At $t=0$, every $\Phi_H(0)=1$, so the positive and negative
coefficient sums are equal.  Separating the two sides and normalizing by
their common sum produces two different probability mixing measures with
the same count probability-generating function, and hence the same MiSo law.
\end{proof}

Theorem~\ref{thm:G-identification-main} is the standard finite-mixture
criterion of \citet{YakowitzSpragins1968}, within the broader identifiability
program of \citet{Teicher1960,Teicher1963}.  It separates two questions that
are sometimes conflated: injectivity of each component kernel, established by
Theorem~\ref{thm:H-identification-main}, and linear independence of the whole
component family, which is required to identify $G$.

\begin{corollary}[Common simplicial dictionary]
\label{cor:dictionary-main}
Fix linearly independent vectors
$\theta_1,\ldots,\theta_K\in\R_+^M\setminus\{0\}$ with $K\leq M$, write
$\Theta=\{\theta_1,\ldots,\theta_K\}$, and consider
\begin{equation}
  \mathcal H_\Theta
  =\left\{
    H_\alpha=\sum_{k=1}^{K}\alpha_k\delta_{\theta_k}:
    \alpha\in\R_+^K\setminus\{0\}
  \right\},
  \label{eq:H-theta-main}
\end{equation}
where a zero coordinate means that the corresponding atom is omitted.  Then
every finite MiSo mixing measure $G$ on $\mathcal H_\Theta$ is identifiable.
\end{corollary}

\begin{proof}
For $H_\alpha\in\mathcal H_\Theta$,
\[
  \Phi_{H_\alpha}(t)
  =\prod_{k=1}^{K}
    (1+\langle\theta_k,t\rangle)^{-\alpha_k}.
\]
Because the $\theta_k$ are linearly independent, the map
$t\mapsto(\langle\theta_1,t\rangle,\ldots,
\langle\theta_K,t\rangle)$ has full row rank and maps an open subset of
$(0,1)^M$ onto an open subset of $\R^K$.  Make the coordinate change
\[
  y_k=\log(1+\langle\theta_k,t\rangle).
\]
A finite linear relation among transforms associated with distinct
$\alpha^{(j)}$ becomes
\[
  \sum_jc_j\exp\{-\langle\alpha^{(j)},y\rangle\}=0
\]
on an open set.  Choose a direction $v$ for which the finitely many values
$\langle\alpha^{(j)},v\rangle$ are distinct and restrict the identity to a
short real line segment parallel to $v$.  Distinct univariate exponentials
are linearly independent, so every $c_j=0$.  The result follows from
Theorem~\ref{thm:G-identification-main}.
\end{proof}

Corollary~\ref{cor:dictionary-main} allows a motif to omit a global ray by
setting its corresponding shape to zero.  It requires the scaled direction
$\theta_k$ to be shared across motifs, as occurs when selecting factor row
$F_k$ also selects a factor-specific rate $\beta_k$.  The result also permits a
fixed over-specified dictionary with $K'>K$ provided all $K'$ scaled
directions remain linearly independent, so necessarily $K'\leq M$.  The true
model is then embedded by assigning zero shape to the extra directions.
However, the corollary conditions on a fixed dictionary; it does not by itself
establish joint identifiability when the dictionary is learned.

\section{Experiments Needed}
\label{sec:experiments}

The empirical goal is not only to show good reconstruction error. The paper
needs to show that MiSo recovers meaningful factor-usage motifs, clusters
observations by motif, learns factors when $F$ is unknown, and quantifies
uncertainty in settings where hard sparse supports are unstable.

\subsection{Baselines}

The main baselines should be:
\begin{enumerate}
  \item Poisson NMF followed by clustering of normalized loading scores;
  \item row-wise Poisson SuSiE with the frozen NMF dictionary, followed by
  clustering of posterior expected loading scores;
  \item sparse NMF baselines with tuned sparsity penalties, if available;
  \item ablations of MiSo with fixed versus learned $F$.
\end{enumerate}
The first two baselines are already implemented in the current benchmark code.

\subsection{Metrics}

The simulation tables should include:
\begin{enumerate}
  \item factor recovery, measured by cosine matching between fitted rows of
  $F$ and true rows of $F_0$;
  \item motif clustering accuracy, measured by matching $\arg\max_s\omega_{is}$
  to true groups;
  \item support recovery, measured by converting posterior motif scores to
  predicted active factors;
  \item soft clustering diagnostics, based on group-level averages of
  $\omega_{is}$ rather than only hard labels;
  \item posterior uncertainty summaries for $\bar{\gamma}_{sd}$, including
  entropy, top-two gaps, and targeted pairwise split metrics;
  \item selected motif dimensions $\widehat{D}_s$ from the ARD prior means;
  \item predictive fit, such as held-out Poisson log likelihood or deviance.
\end{enumerate}

\subsection{Core simulation settings}

The current roadmap suggests the following simulation studies.

\paragraph{Correctly specified block motifs.}
Use correct $K$ and nominal $D$, with groups occupying simple single-factor and
multi-factor supports. This verifies that the method can recover factors,
clusters, and supports in a clean setting. Current evidence shows feasibility,
but a larger seed grid is needed before claiming robust superiority over
baselines.

\paragraph{Tree/no-anchor motifs.}
Use motifs such as $(1,2,3)$, $(1,2,4)$, $(1,5,6)$, $(1,5,7)$, $(1,2)$, and
$(1,5)$, so that no group is explained by a single anchor factor. This is the
strongest current positive setting: in a small multi-seed pilot, MiSo has
substantially better support recovery than the baseline pipelines.

\paragraph{Overspecified $K$.}
Generate data with $K$ true factors but fit with extra candidate factor rows,
including near-duplicates of true factors. The key diagnostic is whether
$\bar{\gamma}_{sd}$ splits posterior mass across duplicate rows when they are
statistically interchangeable. Global entropy is not sufficient; the paper
should report pair-specific posterior split metrics for the duplicated factors.

\paragraph{Correlated factors.}
Generate correlated factor rows at controlled strengths. This tests whether
MiSo can express targeted ambiguity between correlated factors. The current
evidence suggests that raw entropy can be misleading because weak dimensions may
become globally diffuse. The final paper version should sweep correlation strength and
report targeted pair uncertainty, not only mean entropy.

\paragraph{Overspecified $D$.}
Fit each motif with a nominal dimension larger than the true motif dimension.
The target behavior is that extra dimensions either have small ARD prior mean
$\alpha_{sd}^0/\beta_{sd}^0$ or small posterior loading mass. The current
initialization uses the top $D$ center factors as distinct alignment anchors
and averages the aligned row-wise Poisson SuSiE posteriors. A full paper needs
multi-seed threshold sweeps for $\widehat{D}_s$ and comparisons against the
true motif dimensions after matching fitted motifs to groups.

\subsection{Real data}

The real-data analysis should demonstrate three advantages of the MiSo
representation beyond reconstruction quality. First, it should identify
recurring motif types: scientifically meaningful combinations of factors that
are reused across observations. A small atlas of these motifs, together with
their responsibilities, should provide a more concise summary than a typical
admixture bar plot with many separately colored factors. Second, it should show
that a motif represents a low-dimensional submanifold rather than merely a
cluster. In the single-cell application, cells assigned to the same motif
should populate its posterior dimension coordinates continuously, with related
cell states, marker expression, and---where available---spatial position varying
smoothly along a trajectory consistent with differentiation. Third, the
analysis should recover structure among the factors themselves. Factor co-use
across motifs induces a graph in which broadly reused factors may act as hubs
and more specialized factors as leaves; these relationships should be checked
against external biological interpretation.

These are goals for the final analysis rather than claims established by the
current experiment. The identifiability theory remains to be completed, and
the post-processing needed when $K$, $S$, and $D$ are over-specified still
requires further study. In particular, the current Wasserstein dendrogram and
merging rule are exploratory and should not yet be presented as the final
procedure.

\paragraph{A concise empirical demonstration.}
The main real-data result can be organized as one multi-panel figure. The first
panel should show the conventional NMF loading bar plot and its loss of
readability when $K$ is moderately large. The second should replace it with a
small motif atlas, displaying the factors and effective dimension used by each
recurring motif together with its prevalence. The third should highlight one
well-supported two- or three-dimensional motif and plot its cells in the
corresponding posterior dimension coordinates. Cell annotations should be used
only after fitting: related states should occupy a connected trajectory, and
independently chosen marker genes should vary smoothly along the resulting
MiSo coordinate. The fourth panel should show the factor co-use graph, with a
small number of biologically interpretable hub, branch, and leaf factors.

A companion table should report only the quantities needed to support this
picture: held-out deviance relative to Poisson NMF, the average number of active
factors per cell, the effective dimensions of the highlighted motifs,
stability across random seeds, and external agreement with cell annotations
or an independently estimated pseudotime. Thus the intended claim is not that
MiSo improves reconstruction at all costs, but that it preserves predictive fit
while replacing a high-dimensional loading display by recurring motifs,
within-motif coordinates, and a factor-relation graph. The primary display
should use fitted motifs before any uncertain merging step; alternative
post-processing rules can be reported separately as a sensitivity analysis.

\subsection{Ablations and robustness checks}

The paper should include ablations for:
\begin{enumerate}
  \item fixed versus learned $F$;
  \item sensitivity to the number of motifs $S$;
  \item sensitivity to nominal $D$ and the ARD threshold $\tau$;
  \item sensitivity to duplicate-factor and correlated-factor strength;
  \item stability across random seeds.
\end{enumerate}
These experiments are needed to distinguish the proposed statistical model from
an optimization artifact.

\section{Discussion}

MiSo reframes sparse NMF as a motif discovery problem. Instead of estimating
sparse loadings independently for each observation, it pools repeated sparsity
patterns through latent motifs. The posterior motif probabilities
$\bar{\gamma}_{sdk}$ provide an interpretable uncertainty summary over the
factor identities defining each motif, while the responsibilities $\omega_{is}$
cluster observations by motif. The empirical-Bayes prior means
$\alpha_{sd}^0/\beta_{sd}^0$ provide a natural ARD-style diagnostic for the
effective dimension of each motif.

The fitting pipeline separates the roles of its components: Poisson NMF
estimates the initial dictionary, fixed-dictionary row-wise Poisson SuSiE
quantifies observation-level factor uncertainty, and MiSo pools that
information into recurring motifs. This modular separation avoids interleaving
dictionary estimation with observation-level support learning.

The current evidence supports the promise of this approach, especially in
tree/no-anchor simulations and in overfitted-factor uncertainty diagnostics.
The main open issues are stable dimension selection, broader multi-seed
benchmarks, model-selection guidance for $K$, $S$, and $D$, and a real-data
example with scientifically interpretable motifs.

\clearpage
\appendix

\section{Threshold-and-recycle initialization}
\label{app:threshold-recycle}

For reproducibility, we record an earlier initialization rule that is retained
in the software for historical ablations. Using the motif center $c_s$ and
factor ordering $r_s$ defined in the main text, it first forms the ordered set
\[
  A_s(\tau_{\mathrm{init}})
  =
  (a_{s1},\ldots,a_{sq_s})
  =
  \{r_s(j):c_{s,r_s(j)}\geq\tau_{\mathrm{init}}\}.
\]
If this set is empty, it is replaced by $A_s=(r_s(1))$. The rule then fills the
$D$ dimensions cyclically,
\[
  \kappa_{sd}^{\mathrm{recycle}}
  =
  a_{s,\,1+((d-1)\bmod q_s)},
  \qquad d=1,\ldots,D.
\]
When $q_s=1$, every dimension is assigned the same factor. Conditional on motif
$s$, the resulting factor-space loading is
\[
  L_i
  =
  \sum_{d=1}^D\lambda_{id}e_{a_{s1}}
  =
  \left(\sum_{d=1}^D\lambda_{id}\right)e_{a_{s1}},
\]
which lies on a one-dimensional ray regardless of the nominal value of $D$.
More generally, recycling constrains the initialization to at most $q_s$
distinct directions. This can create an artificially low-dimensional starting
point when only a few center weights exceed $\tau_{\mathrm{init}}$.

\bibliographystyle{plainnat}
\bibliography{sparsity}

\end{document}
